% Resum del treball (200-300 paraules aprox.).
%
% Ha de respondre, en poques frases cadascuna, a:
%   - Quin problema es resol i per què és rellevant?
%   - Quin objectiu té el treball?
%   - Quin mètode o enfocament s'ha seguit?
%   - Quins resultats principals s'han obtingut?
%   - Quina és la conclusió clau?


En l'àmbit de l'\textit{Edge Computing}, els dispositius perimetrals solen disposar de recursos molt limitats. El desplegament d'orquestradors de contenidors moderns genera una reserva permanent de memòria i CPU, fet que redueix el marge disponible per a processar càrregues de treball exigents, com ara models d'Intel·ligència Artificial.

L'objectiu principal d'aquest treball és analitzar l'impacte sobre el rendiment que introdueix Kubernetes en entorns de maquinari limitat. Per fer-ho, es compara aquest consum de recursos amb el d'una arquitectura nativa basada en serveis de Systemd, amb la finalitat de determinar quina solució resulta més eficient en cadascun dels escenaris analitzats.

Pel que fa a la metodologia, s'han dissenyat i executat proves d'estrès automatitzades per a ambdues arquitectures utilitzant protocols de comunicació com HTTP-JSON, gRPC-Protobuf i ZeroMQ. L'avaluació s'ha centrat a mesurar el consum de memòria RAM en repòs, la degradació del rendiment de xarxa (\textit{throughput}) i la capacitat d'autorecuperació del sistema. Per a aquesta última, s'ha provocat la caiguda forçada de processos mitjançant el senyal \texttt{SIGKILL} (o \texttt{kill -9}) per tal de calcular el temps mitjà de recuperació (\textit{Mean Time To Repair}, MTTR).

Els resultats indiquen que Kubernetes penalitza notablement els recursos disponibles del node perimetral. L'orquestrador incrementa l'ús base de memòria en un 16~\% (aproximadament 850~MB enfront dels 730~MB de Systemd) per la necessitat de mantenir actiu el seu pla de control (\textit{control plane}). Així mateix, la seva capa de xarxa virtual (\textit{Software-Defined Networking}, SDN) redueix la capacitat de processament en un 30~\% \textemdash{}rebaixant el flux a unes 6.500 imatges per segon enfront de les 9.300 assolides de manera nativa\textemdash{}. Pel que fa a l'autorecuperació, la gestió directa a nivell de sistema operatiu de Systemd aconsegueix reiniciar els serveis afectats tres vegades més ràpid (\textasciitilde 490~ms) que el model distribuït de Kubernetes (\textasciitilde 1600~ms).

A tall de conclusió, Systemd es postula com la solució òptima per a nodes \textit{Edge} aïllats on el rendiment brut i la baixa latència són crítics. El consum addicional de recursos que imposa Kubernetes només es justifica en entorns on es prioritza la simplicitat de gestió, com ara flotes distribuïdes que requereixen desplegaments declaratius i una orquestració centralitzada.